\section{Related Work}
\label{sec:related_work}

\subsection{Digital BE-FAST Screening and Neurological Task Automation}

BE-FAST organizes balance, visual, facial, arm, and speech warning signs into a
rapid public screening procedure and emphasizes timely action
\cite{aroor2017befast}, but correct task execution and interpretation remain
user dependent. Digital studies have automated facial video and speech
\cite{cai2022deepstroke}, landmark sequences from neurological examinations
\cite{lee2022skeleton}, and multimodal limb, face, and speech evidence
\cite{ou2025multimodal,hoang2026digitalfast}. Mobile modified-NIHSS monitoring
has also proved feasible after stroke \cite{noch2023destroke}. This work
establishes task-level feasibility, yet remains distributed across individual
signs, NIHSS subsets, hospital classification, and post-stroke monitoring; a
home workflow must additionally separate valid negative findings from unusable
attempts.

\subsection{Camera- and Audio-Based Neurological Measurement}

Camera landmarks provide compact, inspectable alternatives to end-to-end video
classification. Skeleton sequences represent neurological movements
\cite{lee2022skeleton}; Kinect sway supports controlled balance measurement
\cite{yeung2014kinect}; and face or iris landmarks capture smile motion and
relative eye position \cite{bazarevsky2020mediapipeiris}. Interpretation remains
device and task specific: monocular trunk motion is not force-plate centre of
pressure, and low-rate gaze endpoints are not clinical imaging or
video-oculography \cite{singer2006gaze,kudo2021eyemovement}.

Speech complements vision because post-stroke dysarthria affects timing,
phonation, articulation, and resonance \cite{decock2021acute,kim2011acoustic}.
Mandarin studies report formant and vowel-space changes
\cite{mou2018mandarin,ge2021vowels}, while the Mandarin Dysarthria Speech Corpus
(MDSC) supports learned representations \cite{gao2024mdsc}. Evidence from controlled vowels,
corpus-level labels, and supervised recordings nonetheless requires validation
before transfer to unsupervised home audio and acute screening.

\subsection{Quality-Aware Multimodal Edge Screening}

Multimodal fusion broadens symptom coverage: DeepStroke combined face and
speech, while asymmetric attention reduced noisy-audio influence
\cite{cai2022deepstroke,ou2025multimodal}. For unsupervised screening, however,
normal observations, poor signal, incomplete tasks, and missing modalities must
remain distinguishable. Edge processing reduces network dependence and raw-data
transmission but adds latency, memory, power, and thermal constraints;
RhythmEdge demonstrates relevant Raspberry Pi feasibility and profiling
\cite{hasan2022rhythmedge}. Few stroke-screening designs jointly address these
quality, missingness, and scheduling requirements.

\paragraph{Research gap.}
Prior work establishes component feasibility, but integration is incomplete. PiBE-FAST
targets a complete edge-based BE-FAST workflow, including balance and ocular
tasks, valid-negative versus insufficient-result handling, modality-specific
quality gates, explicit missingness, stage-aware local inference, and separation
of research fusion from user-facing urgency rules.
